\documentclass[11pt,a4paper]{article}
\usepackage[utf8]{inputenc}
\usepackage[T1]{fontenc}
\usepackage[spanish,es-noquoting]{babel}
\usepackage[margin=2.5cm]{geometry}
\usepackage{booktabs}
\usepackage{array}
\usepackage{xcolor}
\usepackage{colortbl}
\usepackage{enumitem}
\usepackage{fancyhdr}
\usepackage{titlesec}
\usepackage[hidelinks]{hyperref}
\usepackage{tcolorbox}

\definecolor{acc}{HTML}{0E6E65}
\definecolor{cop}{HTML}{A85A32}
\definecolor{ln}{HTML}{D5E2DF}
\definecolor{ok}{HTML}{2E7D5B}
\definecolor{no}{HTML}{B23A48}

\titleformat{\section}{\normalfont\large\bfseries\color{acc}}{\thesection.}{0.5em}{}
\titleformat{\subsection}{\normalfont\normalsize\bfseries}{\thesubsection}{0.5em}{}

\pagestyle{fancy}\fancyhf{}
\lhead{\small\color{acc}swarm-abm vs.\ sim2-agricultores}
\rhead{\small\thepage}
\renewcommand{\headrulewidth}{0.4pt}

\newcommand{\si}{\textcolor{ok}{\textbf{sí}}}
\newcommand{\nno}{\textcolor{no}{\textbf{no}}}
\newcommand{\code}[1]{\texttt{\small #1}}

\begin{document}

\begin{center}
{\LARGE\bfseries\color{acc} Idoneidad de un motor C++ existente para\\[2pt]
escalar el modelo de depredación de ovejas}\\[10pt]
{\large Análisis técnico comparativo}\\[14pt]
\begin{tabular}{rl}
\textbf{Preparado para:} & Prof.\ Mauricio Marín \\
 & Manuel Manríquez (autor de \code{sim2-agricultores}) \\[3pt]
\textbf{Preparado por:} & Francisco Parra \\[3pt]
\textbf{Fecha:} & 13 de julio de 2026 \\[3pt]
\textbf{Objeto:} & Contrastar el motor C++ propuesto contra \code{swarm-abm} \\
 & como vehículo para escalar el modelo SIGRID (depredación de ganado). \\
\end{tabular}
\end{center}

\vspace{4pt}
\hrule
\vspace{10pt}

\begin{tcolorbox}[colback=acc!6,colframe=acc,boxrule=0.6pt,arc=2pt]
\textbf{Resumen ejecutivo.} Para \emph{escalar} el problema de las ovejas
depredadas se requiere calibrar el modelo y hacerle análisis de sensibilidad de
forma auditable, y ambas cosas descansan en la \textbf{reproducibilidad}
(mismo input $\rightarrow$ mismo output). Una prueba en vivo muestra que el
código C++ propuesto, \emph{tal como está hoy}, \textbf{no reproduce sus propias
corridas}: dos ejecuciones con configuración idéntica dan resultados distintos.
El mismo experimento sobre \code{swarm-abm} es reproducible bit a bit. Esto no
significa que el C++ sea ``inutilizable'': la falla es arreglable en principio.
El punto es de \textbf{esfuerzo y riesgo}: usarlo para esta tarea implica
reconstruir en C++ el motor espacial, el diseño de experimentos y la garantía de
reproducibilidad que \code{swarm-abm} ya provee, tiene auditado y publicado.
\end{tcolorbox}

\section{Qué es el código C++ propuesto}

\code{sim2-agricultores} (\textasciitilde5.160 líneas de C++) es un simulador de
\textbf{mercados agrícolas} —agricultores, feriantes, consumidores, mercado
mayorista, terrenos y productos, con eventos de sequía, helada y ola de calor— y
es una pieza sólida en su dominio. Su arquitectura es de \textbf{simulación por
eventos discretos} (DES): una \emph{future event list}, manejadores de eventos y
colas de mensajes.

El modelo de ovejas (SIGRID) es de naturaleza distinta: \textbf{espacial y por
pasos} (grilla, dinámica depredador--presa evaluada por \emph{tick}). De ahí dos
observaciones de fondo, previas a la reproducibilidad:

\begin{itemize}[leftmargin=1.4em,itemsep=2pt]
  \item El código C++ \textbf{no contiene primitivas espaciales} (grilla,
  vecindades). Llevar SIGRID ahí obliga a construir esa infraestructura desde
  cero.
  \item El \textbf{paradigma no calza}: un DES orientado a transacciones de
  mercado no expresa naturalmente la dinámica espacial sincrónica de depredación.
\end{itemize}

\section{Prueba en vivo: el C++ propuesto no es reproducible}

Se compiló y ejecutó el simulador C++ \textbf{dos veces con configuración
idéntica} (\code{sim\_config.json} sin cambios), en un contenedor autocontenido
(PostgreSQL~16 + g++~13 + libpqxx). Se comparó el vector completo de resultados
agregados mediante hash MD5:

\begin{tcolorbox}[colback=black!3,colframe=ln,boxrule=0.5pt,arc=2pt,fontupper=\ttfamily\small]
Corrida 1 — exit 0 — 1260 filas — hash \textcolor{cop}{089cdd94c7695f77\dots} \\
Corrida 2 — exit 0 — 1260 filas — hash \textcolor{cop}{9f344b0f1e960c3d\dots}\ \ $\leftarrow$ distinto \\[4pt]
65 de 1260 filas (misma clave) tienen VALOR distinto entre las dos corridas. \\
Ej.: ``COMPRA DE FERIANTE A AGRICULTOR'', t=16, prod=18 $\rightarrow$ 36.360 vs 28.620 ($-21\%$) \\[4pt]
\textbf{VEREDICTO: NO REPRODUCIBLE.}
\end{tcolorbox}

\noindent\textbf{Causa.} Cada fuente de aleatoriedad del código se siembra con
\code{std::mt19937(std::random\_device\{\}())} —desde la entropía del sistema
operativo, sin semilla fija en ninguna parte del API (aparece en al menos 8
archivos)—. Con el código actual no hay forma de forzar dos corridas idénticas.

\medskip
\noindent\textit{Matiz honesto.} Para estadísticas agregadas de mercado por Monte
Carlo (correr $N$ veces y promediar), esta decisión es legítima. Pero es una
limitación para calibración y análisis de sensibilidad —justamente lo que
``escalar el problema'' requiere— y para depurar (un resultado anómalo no se
puede reproducir).

\section{Prueba en vivo: swarm-abm sí es reproducible}

El mismo experimento sobre el modelo de ovejas real en \code{swarm-abm} (dos
corridas con la misma semilla):

\begin{tcolorbox}[colback=black!3,colframe=ln,boxrule=0.5pt,arc=2pt,fontupper=\ttfamily\small]
Corrida 1 — hash de la salida: \textcolor{acc}{faa5b35e733da2cb\dots} \\
Corrida 2 — hash de la salida: \textcolor{acc}{faa5b35e733da2cb\dots}\ \ $\leftarrow$ idéntico \\[4pt]
\textbf{VEREDICTO: REPRODUCIBLE BIT A BIT.}
\end{tcolorbox}

\noindent No es accidente de una corrida: el determinismo de \code{swarm-abm}
está garantizado por construcción (RNG ChaCha8 sembrable, semilla por-agente,
ejecución paralela bit-idéntica a la secuencial, identidad entre x86-64 y
WebAssembly) y \textbf{verificado en integración continua} en cada cambio.

\section{Matriz de idoneidad para escalar el modelo de ovejas}

\renewcommand{\arraystretch}{1.3}
\begin{center}
\begin{tabular}{>{\raggedright\arraybackslash}p{7.4cm} c c}
\toprule
\textbf{Lo que ``escalar'' exige} & \textbf{sim2-agricultores} & \textbf{swarm-abm} \\
\midrule
Reproducibilidad bit a bit & \nno\ (demostrado) & \si\ (demostrado + CI) \\
Ya implementa el modelo de ovejas & \nno\ (es mercados) & \si\ (SIGRID vs Mesa) \\
Primitivas espaciales (grilla, vecindades) & \nno\ (DES no espacial) & \si\ (tres espacios) \\
Análisis de sensibilidad global nativo & \nno\ (a construir) & \si\ (Sobol/Morris/LHS) \\
Calibración & \nno\ (manual) & \si\ (DE, $\sim$400$\times$ vs Python) \\
Determinismo bajo paralelismo & \nno & \si\ (paralelo $=$ secuencial) \\
Dependencias para correr & PostgreSQL & ninguna ($+$ WASM) \\
Estado del código & tesis, \code{-fpermissive} & crates.io, 3 auditorías, CI \\
\bottomrule
\end{tabular}
\end{center}

\section{Sobre el rendimiento (honesto)}

\textbf{No es posible comparar la velocidad de ambos directamente}: son modelos y
paradigmas distintos, y cronometrar uno contra otro no mediría nada. Lo que sí
está medido con rigor es que \code{swarm-abm} \emph{no sacrifica rendimiento por
ser reproducible}: es 45--184$\times$ más rápido que Mesa (Python) y está a solo
$\sim$1,5--2,6$\times$ del motor de ABM nativo más rápido que existe (krABMaga,
Rust). Es decir, el argumento ``hay que ir a C++ por velocidad'' no se sostiene:
un motor en Rust ya está en la banda de rendimiento nativo \emph{y además} es
reproducible y trae análisis de sensibilidad integrado.

\section{Conclusión}

No se afirma que el código C++ sea inutilizable —sería impreciso—. Hay dos cosas
distintas: (i)~la no-reproducibilidad es un defecto \textbf{arreglable en
principio}, cableando una semilla, aunque hacerlo bien (y mantener determinismo
si se paraleliza) es trabajo de ingeniería no trivial; y (ii)~el obstáculo mayor
es que el código \textbf{no es el modelo ni el paradigma}: usarlo para las ovejas
implica reimplementar el modelo espacial depredador--presa en un framework de
mercados que no lo soporta.

\begin{tcolorbox}[colback=cop!7,colframe=cop,boxrule=0.6pt,arc=2pt]
El encuadre correcto no es ``se puede o no se puede'', sino \textbf{esfuerzo y
riesgo frente a reutilizar lo ya hecho}. Usar este C++ para escalar el problema
significa reconstruir —con menos garantías y semanas de trabajo— el motor
espacial, el diseño de experimentos y la reproducibilidad que \code{swarm-abm} ya
provee, tiene auditado y publicado. El modelo de ovejas ya corre ahí, validado
contra Mesa y reproducible bit a bit, y fue con su análisis de sensibilidad
nativo que se identificó al número de perros como la principal palanca de manejo.
\end{tcolorbox}

\vspace{6pt}
\noindent\footnotesize\textcolor{gray}{La comparación es reproducible: el arnés
de la prueba (Dockerfile + script) está disponible y repite ambos lados en un
comando. swarm-abm: \url{https://github.com/franciscoparrao/swarm-abm} /
\url{https://crates.io/crates/swarm-abm}.}

\end{document}
